\textbf{Transfer to unseen crack counts and sensitivity to graph connectivity.}
\textbf{a}, Equal-case mean absolute error (MAE) for each held-out crack count, $N=9,\ldots,15$, after training and validation on $N\leq8$. All seven models are shown in the frozen protocol order on a common vertical scale. Hollow circles show the five training seeds, with small horizontal offsets solely for visibility; filled markers and lines indicate their arithmetic means. The 264 test cases all originate from source batch 005, so source batch and high crack count are confounded. Case counts for increasing $N$ are 48, 45, 41, 40, 36, 29, and 25.
\textbf{b}, Changes in equal-case MAE when the same trained GNN checkpoint is evaluated on radius and symmetric $k$-nearest-neighbor graph views, relative to its complete-graph result. Thin lines pair the values for one seed; dark ticks show means. All three prespecified GNN variants are reported on a common scale. Radius and $k$ are fixed by the recorded split recipe, with no refitting on the sparse views. A positive change denotes higher error. These connectivity tests reuse the same held-out cases and test representation sensitivity, not new finite-element geometries or mesh transfer.
All plotted errors are percentage points of the dimensionless archived scalar-envelope target. Seed points describe training variability and are not confidence intervals over cases; time points are not pooled as independent observations.
